\chapter{Атрибуция Hodge-уровня фоновому оператору H15}
\label{ch:v7-hodge-level-background-attribution-gate}

\section{Вопрос относительного масштаба}

Полевой Hodge-родитель содержит уровень
\begin{equation}
 \mathfrak m_\mu(Z)
 =[d_Z,d_Z^\dagger]-\mu^2\widehat\Gamma_E.
 \label{eq:v7-background-level-moment}
\end{equation}
Абсолютное значение $\mu$ уже запрещено выводить из одной безразмерной
конечной геометрии. Остаётся более слабый вопрос: может ли оператор
$D_{H_{15}}$ единственным образом определить отношение
$\mu/\|D_{H_{15}}\|$?

После доказанной ортогональности трёх физических кромок типизированный фон
есть прямая сумма трёх двухступенных комплексов
\begin{equation}
 d_{15}=d_u\oplus d_d\oplus d_e,
 \qquad
 K_{15}=[d_{15},d_{15}^\dagger],
 \qquad
 \chi_{15}=\operatorname{diag}(-I_3,I_3).
 \label{eq:v7-background-level-typed-complex}
\end{equation}
Пусть $\Pi_a$ выделяет пару концов ребра $a\in\{u,d,e\}$. Тогда три
положительные линейные координаты момента равны
\begin{equation}
 k_a(K_{15})
 =\frac12\Tr\!\left(\chi_{15}\Pi_aK_{15}\Pi_a\right)
 =\Tr(Y_a^\dagger Y_a)\ge0.
 \label{eq:v7-background-level-edge-energies}
\end{equation}

\section{Классификация минимальных эквивариантных отображений}

Целевой уровень должен лежать на уже выведенной операторной прямой
$\mathbb R\widehat\Gamma_E$. Уже минимальный класс линейных
самосопряжённых отображений, использующий три типизированных фоновых момента,
содержит семейство
\begin{equation}
 \Psi_c(K_{15})
 =\left(c_uk_u+c_dk_d+c_ek_e\right)\widehat\Gamma_E,
 \qquad c=(c_u,c_d,c_e)\in\mathbb R^3.
 \label{eq:v7-background-level-map-family}
\end{equation}
Для положительного уровня на положительном конусе необходимо и достаточно
$c_a\ge0$.

Эти три коэффициента не отождествляются симметрией. Рёбра $u,d,e$ несут
разные точные калибровочные типы, поэтому орбиты типосохраняющей группы
равны
\begin{equation}
 \{u\},\qquad\{d\},\qquad\{e\}.
 \label{eq:v7-background-level-exact-orbits}
\end{equation}
Следовательно, минимальное пространство допустимых отображений содержит
трёхмерное подпространство. Даже после забывания гиперзаряда остаются две
орбиты
\begin{equation}
 \{u,d\},\qquad\{e\},
 \label{eq:v7-background-level-coarse-orbits}
\end{equation}
и двумерное пространство отображений. После условия нормировки
$c_u+c_d+c_e=1$ положительное точное семейство остаётся двумерным
симплексом, а грубое $u\leftrightarrow d$-симметричное семейство ---
одномерным интервалом.

\begin{theorem}[Неединственность фоновой атрибуции]
Gauge-, Real- и градуировочная эквивариантность не выделяют единственного
линейного отображения $K_{15}\mapsto\mathbb R\widehat\Gamma_E$. Уже в
минимальном классе допустимо трёхмерное
семейство~\eqref{eq:v7-background-level-map-family};
одна положительная нормировка уменьшает его только до двумерного симплекса.
Единственность потребовала бы дополнительной перестановочной симметрии трёх
физически неэквивалентных рёбер.
\end{theorem}

\section{Две независимые неоднозначности}

В семейно-слепом единичном фоне
\begin{equation}
 (k_u,k_d,k_e)=(1,1,1),
 \qquad
 \Tr D_{H_{15}}^2=6.
 \label{eq:v7-background-level-unit-background}
\end{equation}
Даже равновесовой луч не получает единственной численной нормировки:
\begin{equation}
 \frac1{9}\Tr D_{H_{15}}^2=\frac23,
 \qquad
 \frac1{5}\Tr_{\rm active}D_{H_{15}}^2=\frac65,
 \qquad
 \frac1{2|E_0|}\Tr D_{H_{15}}^2=1.
 \label{eq:v7-background-level-trace-ambiguity}
\end{equation}
Это свобода общей меры между физическим и полевым носителями.

В общем семейно-неравном фоне возникает дополнительная свобода направления.
Например, при
\begin{equation}
 (k_u,k_d,k_e)=(1,4,9)
 \label{eq:v7-background-level-unequal-example}
\end{equation}
три строго положительных нормированных выбора
$c=(1,1,1)/3$, $c=(1,1,2)/4$ и $c=(1/2,1/3,1/6)$ дают разные уровни.
Значит, совпадение на точке $(1,1,1)$ не продолжается канонически на
физические неравные юкавские матрицы.

Условие изометричности также не спасает единственность: линейный функционал
из трёхмерного пространства в одну прямую задаётся единичным вектором
$c\in S^2$. Положительность оставляет непрерывный сектор этой сферы.
Следосохранение фиксирует одну линейную комбинацию, но не оставшиеся две.

\section{Литературная граница}

Квадрат нормы отображения момента задаёт корректную динамику после выбора
центрального уровня, но сам уровень является параметром устойчивости
колчанного quotient \cite{v7-background-level-harada}. Спектральное действие
аналогично содержит масштаб как исходное данное; замена его дилатоном делает
масштаб динамическим полем, но не выводит единственное вакуумное значение из
конечного графа \cite{v7-background-level-dilaton}. Поэтому найденная
неединственность является ожидаемой структурной границей, а не дефектом
конкретной координатизации.

\begin{nogocriterion}[Закрытие внутренней атрибуции уровня]
Запрещено выбирать $c_u=c_d=c_e$, одну из нормировок
\eqref{eq:v7-background-level-trace-ambiguity} или единичный вектор $c$ по
желаемому масштабу. Фон $H_{15}$ даёт ненулевые масштабные наблюдаемые, но
не единственное отображение между его носителем и
$\widehat\Gamma_E$. Поэтому относительный уровень $\mu/\|D_{H_{15}}\|$ в
текущей конечной архитектуре не выведен.
\end{nogocriterion}

\begin{protocol}
\begin{enumerate}
 \item независимых типизированных фоновых моментов: \textbf{$3$};
 \item ранг минимального пространства $\Psi_c$: \textbf{$3$};
 \item размерность после грубой симметрии $u\leftrightarrow d$:
 \textbf{$2$};
 \item положительное нормированное точное семейство: \textbf{двумерное};
 \item три естественные следовые нормировки совпадают: \textbf{нет};
 \item общий неравный фон устраняет свободу весов: \textbf{нет};
 \item Real- и градуировочная совместимость: \textbf{проверена};
 \item единственное отношение $\mu/\|D_{H_{15}}\|$: \textbf{не получено};
 \item абсолютная шкала: \textbf{не получена};
 \item итог: \textbf{no-go внутренней атрибуции Hodge-уровня}.
\end{enumerate}
\end{protocol}

Машинный сертификат сохранён в
\path{s2t_v7_hodge_level_background_attribution_gate_results.json}.

\begin{thebibliography}{9}
\bibitem{v7-background-level-harada}
M.~Harada, G.~Wilkin,
\emph{Morse Theory of the Moment Map for Representations of Quivers},
Geometriae Dedicata 150 (2011), 307--353; arXiv:0807.4734.
\bibitem{v7-background-level-dilaton}
A.~H.~Chamseddine, A.~Connes,
\emph{Scale Invariance in the Spectral Action},
Journal of Mathematical Physics 47 (2006), 063504;
arXiv:hep-th/0512169.
\end{thebibliography}